% Part: model-theory
% Chapter: lindstrom
% Section: abstract-logics

\documentclass[../../../include/open-logic-section]{subfiles}

\begin{document}

\olfileid[fr]{mod}{lin}{alg}
\olsection{Logiques abstraites}

\begin{defn}
Une \emph{logique abstraite} est une paire
$\tuple{L,\models_L}$, où~$L$ associe
à chaque !!{language}~$\Lang{L}$ un
ensemble $L(\Lang{L})$ d'énoncés,
et où~$\models_L$ est une relation
de satisfaction entre les
!!{structure}s de~$\Lang{L}$ et les
éléments de~$L(\Lang{L})$.
En particulier, $\tuple{F,\models}$
désigne la logique ordinaire du premier
ordre : $F(\Lang{L})$ est l'ensemble
des !!{sentence}s du premier ordre
construites dans~$\Lang{L}$, et
$\models$ est leur relation habituelle
de satisfaction.
\end{defn}

Nous conservons la notion de
!!{structure} d'un !!{language} du
premier ordre. En revanche, nous
n'imposons ni que les énoncés de
$L(\Lang{L})$ soient construits
de la manière habituelle à partir
des symboles de~$\Lang{L}$, ni
que leur satisfaction soit définie
récursivement comme au premier ordre.
Cette définition générale admet,
par exemple, des conjonctions ou
disjonctions infinies, des
quantificateurs autres que~$\lexists$
et~$\lforall$ (comme « il existe
une infinité de~$x$ tels que… »),
ou encore des préfixes infinis de
quantificateurs. Pour distinguer
ces énoncés des formules ordinaires,
nous les noterons $!E,!F,\dots$;
les lettres $!A,!B,\dots$ seront
réservées aux !!{formula}s du
premier ordre.

\begin{defn}
Notons $\Mod(L){!E}$ la classe
$\Setabs{\Struct{M}}{\Struct{M}\models_L !E}$.
Lorsque le langage doit être précisé,
écrivons $\Mod[L](L){!E}$.
Deux !!{structure}s $\Struct{M}$ et
$\Struct{N}$ de~$\Lang{L}$ sont
\emph{élémentairement équivalentes
dans} $\tuple{L,\models_L}$,
ce que l'on note
$\Struct{M}\elemequiv[L]\Struct{N}$,
si les mêmes énoncés de
$L(\Lang{L})$ sont vrais dans
l'une et dans l'autre.
\end{defn}

\begin{defn}
Une logique abstraite
$\tuple{L,\models_L}$ est dite
\emph{normale} si elle possède
les propriétés suivantes :
\begin{enumerate}
\item (\emph{Monotonie de~$L$})
  Si $\Lang{L}\subseteq\Lang{L}'$,
  alors $L(\Lang{L})\subseteq
  L(\Lang{L}')$; pour un énoncé
  du plus petit langage, l'expansion
  d'une structure aux symboles
  supplémentaires ne change pas
  sa vérité.\footnote{Complément éditorial : la source énonce ici l'inclusion des ensembles d'énoncés. Nous explicitons aussi l'invariance de leur vérité lorsque l'on ajoute des symboles inutilisés.}
\item (\emph{Propriété d'expansion,
  ou dépendance finie}) Pour tout
  $!E\in L(\Lang{L})$, il existe
  un sous-langage \emph{fini}
  $\Lang{L}'\subseteq\Lang{L}$
  et un énoncé $!E'\in L(\Lang{L}')$
  tels que, pour toute
  $\Lang{L}$-structure $\Struct{M}$,
  \[
    \Struct{M}\models_L !E
    \quad\text{si et seulement si}\quad
    \Struct{M}|_{\Lang{L}'}\models_L !E'.
  \]
  Ainsi, deux
  structures ayant le même réduit
  à~$\Lang{L}'$ satisfont toutes
  deux~$!E$ ou aucune ne le satisfait.\footnote{Complément éditorial : outre la dépendance de la vérité à un nombre fini de symboles, nous demandons un représentant dans ce sous-langage fini. La source ne formule que la dépendance, mais utilise ensuite un tel représentant dans les preuves.}
\item (\emph{Invariance par
  isomorphisme}) Si
  $\Struct{M}\models_L !E$ et
  $\Struct{M}\simeq\Struct{N}$,
  alors $\Struct{N}\models_L !E$.
\item (\emph{Invariance par
  renommage}) Un renommage bijectif
  des symboles, préservant leur
  nature et leur arité, transporte
  les énoncés et leurs modèles :
  si $\Lang{L}'$ est le langage
  renommé, $!E'$ le renommage de
  $!E$ et $\Struct{M}'$ la structure
  correspondante, alors
  $\Struct{M}\models_L !E$ si et
  seulement si
  $\Struct{M}'\models_L !E'$.
\item (\emph{Clôture booléenne})
  Pour tout $!E\in L(\Lang{L})$,
  il existe un énoncé~$!F$ du même
  langage tel que
  $\Struct{M}\models_L !F$ si et
  seulement si
  $\Struct{M}\not\models_L !E$.
  Pour $!E$ et~$!F$, il existe
  aussi un énoncé~$!G$ tel que
  $\Mod(L){!G}=\Mod(L){!E}\cap
  \Mod(L){!F}$. La logique représente
  de même les énoncés atomiques
  clos du premier ordre et les autres
  connecteurs booléens.
\item (\emph{Propriété des
  quantificateurs}) Soient une
  constante $c\in\Lang{L}$,
  $!E\in L(\Lang{L})$ et
  $\Lang{L}'=\Lang{L}\setminus\{c\}$.
  Il existe un énoncé
  $!F\in L(\Lang{L}')$ tel que
  \[
  \Mod[L'](L){!F}=
    \Setabs{\Struct{M}}{
    \Expan{M}{a}\in\Mod[L](L){!E}
    \text{ pour un }a\in\Domain{M}},
  \]
  où $\Expan{M}{a}$ est l'expansion
  de~$\Struct{M}$ qui interprète~$c$
  par~$a$.\footnote{Correction éditoriale : l'énoncé $!F$ appartient au langage réduit $\Lang{L}'$, et non au langage initial indiqué par la source, puisque ses modèles sont des $\Lang{L}'$-structures.}
\item (\emph{Propriété de
  relativisation}) Soient
  $!E\in L(\Lang{L})$ et des
  symboles nouveaux $R,c_1,\dots,c_n$,
  où~$R$ a $n+1$ places. Il existe
  un énoncé
  $!F\in L(\Lang{L}\cup
  \{R,c_1,\dots,c_n\})$, appelé
  \emph{relativisation} de~$!E$
  à $\Atom{R}{x,c_1,\dots,c_n}$,
  tel que, pour toute expansion
  $\Expan{M}{X,b_1,\dots,b_n}$
  où
  \[
    D=\Setabs{a\in\Domain{M}}{
      (a,b_1,\dots,b_n)\in X}
  \]
  est non vide et contient les valeurs de toutes les constantes de~$\Lang{L}$ ; il définit alors le domaine d'une
  sous-structure~$\Struct{N}$
  de~$\Struct{M}$, on ait
  \[
    \Expan{M}{X,b_1,\dots,b_n}
      \models_L !F
    \quad\text{si et seulement si}\quad
    \Struct{N}\models_L !E.
  \]
  Ici $R$ est interprété par
  $X\subseteq\Domain{M}^{n+1}$
  et $c_i$ par~$b_i$; la
  sous-structure induite est celle
  de l'\olref[bas][sub]{rem:substructure}.\footnote{Correction éditoriale : le prédicat nouveau $R$ est interprété par $X$ dans l'expansion ; il n'a pas d'interprétation dans la structure de base $\Struct{M}$, contrairement à la notation de la source.}
\end{enumerate}
\end{defn}

\begin{defn}
Soient deux logiques abstraites
$\tuple{L_1,\models_{L_1}}$ et
$\tuple{L_2,\models_{L_2}}$.
La seconde est \emph{au moins aussi
expressive} que la première, ce
que l'on note
$\tuple{L_1,\models_{L_1}}
\leq\tuple{L_2,\models_{L_2}}$,
si, pour tout !!{language}~$\Lang{L}$
et tout $!E\in L_1(\Lang{L})$,
il existe $!F\in L_2(\Lang{L})$
tel que
$\Mod[L](L_1){!E}=
\Mod[L](L_2){!F}$.
Les deux logiques sont
\emph{équivalentes} si chacune
est au moins aussi expressive
que l'autre.
\end{defn}

\begin{rem}
La logique du premier ordre
$\tuple{F,\models}$ est normale.
Les propriétés ci-dessus sont
pour l'essentiel immédiates dans
ce cas. La dépendance finie tient
au fait que la vérité d'un énoncé
ne dépend que de l'interprétation
de ses symboles. Pour relativiser
$!E$ à $\Atom{R}{x,c_1,\dots,c_n}$,
on relativise récursivement ses
sous-formules : une quantification
$\lforall[x][!F]$ devient
$\lforall[x][(\Atom{R}{x,c_1,\dots,c_n}
\lif !F)]$, où~$!F$ a déjà été
relativisée. Les autres cas se
traitent de la même manière ; $\lexists[x][!F]$ devient $\lexists[x][(\Atom{R}{x,c_1,\dots,c_n}\land !F)]$, avec $!F$ déjà relativisée.
Enfin, si $\tuple{L,\models_L}$
est normale, une récurrence sur
les !!{formula}s du premier ordre
montre que
$\tuple{F,\models}\leq
\tuple{L,\models_L}$.
On peut donc, à équivalence
sémantique près, représenter
tout énoncé du premier ordre
dans chaque logique normale.\footnote{Précision éditoriale : les opérations de normalité fournissent des représentants sémantiquement équivalents des énoncés du premier ordre. Leur identité syntaxique avec des éléments de $L(\Lang{L})$ ne découle pas de la seule comparaison d'expressivité ; elle requerrait un choix de représentation.}
\end{rem}

\end{document}
